% 原文 p.93（後半）


\subsectionlink{デコヒーレンスと散逸}
系からのエネルギーの喪失である散逸は、古典的な効果である。系を、環境と呼んでもよい他の系と相互作用させると、両者の間ではしばしばエネルギー交換が起こり、最終的に系は熱平衡に近づく。これが起こる特徴的時間尺度は、通常、その系の緩和時間尺度と呼ばれる。特定の系と環境の相互作用が、系に散逸をもたらす場合、その相互作用の強さは、緩和時間の尺度となる。相互作用が弱くなるほど緩和時間は長くなり、逆も成り立つ。

しかし、散逸がないか無視できるとしても、系と環境の相互作用がない、または無視できると結論してはならない。ある型の相互作用が散逸をもたらさなくても、一般には、純粋な量子効果であるデコヒーレンスを起こす。すなわち環境は一般に、系のエネルギーをまったく吸収せずに、経路情報、または、より一般的にどの状態かという情報を得られる。前の2.10節の模型は、そのような散逸を伴わないデコヒーレンスの、最初の具体例を与えた。したがって、デコヒーレンスは散逸を伴ってもよいが、伴う必要はない。一方、散逸が存在するなら、デコヒーレンスも起こる。

散逸とデコヒーレンスの両方がある場合には、両者の時間尺度は大きく異なるため、通常はかなり容易に区別できる。

% 原文 p.94
\hypertarget{srcpage:94}{}\label{srcpage:94}
多数の理論的・実験的研究が示したように、デコヒーレンス時間は通常、緩和時間より何桁も短い。5.2.4節では、距離$\Delta x$だけ離れた、二つの異なる位置の重ね合わせで記述される、質量を持つ物体について、緩和時間$\tau_r$とデコヒーレンス時間$\tau_d$の比の、おおよその評価が
\begin{equation}\frac{\tau_r}{\tau_d}\sim\left(\frac{\Delta x}{\lambda_{\mathrm{dB}}}\right)^2\tag{2.113}\end{equation}
であることを示す。ここで$\lambda_{\mathrm{dB}}$は、物体の熱的de Broglie波長、
\begin{equation}\lambda_{\mathrm{dB}}=\frac{\hbar}{\sqrt{2mk_{\mathrm B}T}}\tag{2.114}\end{equation}
であり、$m$は物体の質量、$T$は温度、$k_{\mathrm B}$はBoltzmann定数である。式(2.113)は、Zurekが1984年に導出し\cite{ref012}、当時はまだ「脱相関時間」と呼んだデコヒーレンス時間の、最初の一般的評価を与えた。

Zurekはその論文と、その後のPhysics Todayの記事\cite{ref013}で、巨視的水準でのデコヒーレンスと緩和の時間尺度が、通常いかに大きく違うかを示すため、次の数値例を与えた。室温$T=300\,\mathrm K$で、質量$m=1\,\mathrm g$の巨視的物体を考えた。この物体の$\lambda_{\mathrm{dB}}$は微小で、約$10^{-23}\,\mathrm m$である。明らかにこれは、どの意味でも微視的な数である。いま粒子が、中間的、さらには巨視的に区別可能な二位置の、コヒーレントな重ね合わせで記述されるとしよう。すると式(2.113)の比$(\Delta x/\lambda_{\mathrm{dB}})^2$は非常に大きく、この重ね合わせのデコヒーレンス時間$\tau_d$は、散逸、すなわち緩和時間$\tau_r$より圧倒的に短いことが従う。Zurek\cite{ref012,ref013}は、物体が、巨視的な距離$\Delta x=1\,\mathrm{cm}$だけ離れた二つの空間位置の重ね合わせで記述される例を用いた。この場合、比(2.113)は$10^{40}$程度となる。すなわち、その重ね合わせで記述される物体の典型的デコヒーレンス時間は、熱平衡に達する緩和時間より、約40桁も短いのである。

したがって巨視的物体では、環境によるデコヒーレンスに関係する、どの時間尺度での系の力学についても、環境の散逸的影響は、通常まったく無視できる。たとえばボウリングの球に散乱される光子は、球の運動にはほとんど影響しない一方、球の巨視的に識別できる位置を含む重ね合わせ状態には、ほぼ瞬間的なデコヒーレンスをもたらす（図1.3も参照）。

このように、緩和時間とデコヒーレンス時間の決定的な違いは、なぜ巨視的物体が、デコヒーレンスの作用によって実効的に「作り出される」（5.2.5節）、完全にNewton的な軌道を進むと観測されながら、物体の減速など、散逸の現れがほとんどない場合が多いのかを説明する。Joosに由来する例\cite[p.~79]{ref017}を用いれば、木星は何十億年も太陽の周りをNewton的軌道で公転している一方、その運動状態は、表面で散乱される太陽光による散逸的損失、すなわち摩擦に、ほとんど影響されていない。緩和時間とデコヒーレンス時間が実際に同程度なら、明らかに、この事実を説明するのは非常に難しいだろう。

% 原文 p.95
\hypertarget{srcpage:95}{}\label{srcpage:95}
\subsectionlink{デコヒーレンスと古典的雑音}
本書では、一貫して「デコヒーレンス」という語を、何らかの環境との、通常は実際上不可逆な量子もつれの帰結を記述するために限って使う。これは、歴史的に確立した意味、および環境デコヒーレンスの膨大な文献と一致する。2.6節、2.7節で詳しく述べたように、デコヒーレンスの基本機構は、次のように形式化できる。系$\mathcal S$に、環境$\mathcal E$との相互作用の過程で、もつれが最小になるという性質によって定義される、指針状態の組$\{\lvert s_n\rangle\}$があるとする。
\begin{equation}\lvert s_n\rangle\lvert E_0\rangle\longrightarrow\lvert s_n\rangle\lvert E_n\rangle.\tag{2.115}\end{equation}
すると、指針状態$\lvert s_n\rangle$の任意の純粋状態の重ね合わせ$\sum_n c_n\lvert s_n\rangle$は、なじみ深いもつれ状態
\begin{equation}\left(\sum_n c_n\lvert s_n\rangle\right)\lvert E_0\rangle\longrightarrow\sum_n c_n\lvert s_n\rangle\lvert E_n\rangle\tag{2.116}\end{equation}
をもたらす。2.7節で述べたとおり、縮約密度行列の形式では、これは$\{\lvert s_n\rangle\}$基底の非対角項の減衰をもたらす。
\begin{equation}\hat\rho_{\mathcal S}=\sum_{nn'}c_nc_{n'}^*\lvert s_n\rangle\langle s_{n'}\rvert\langle E_{n'}\vert E_n\rangle\longrightarrow\sum_n\lvert c_n\rvert^2\lvert s_n\rangle\langle s_n\rvert.\tag{2.117}\end{equation}
環境の相対状態$\lvert E_n\rangle$が、急速に相互直交、すなわち識別可能になるからである。明らかに、相互作用(2.116)は、$\mathcal S$のどの成分状態にも力学的変化を生まず、唯一の効果は純粋で非局所的な量子もつれである。したがって、デコヒーレンスは、古典的類似物を持たない、明確に量子力学的な効果として理解すべきである。これが、本書で用いるデコヒーレンスの形式と物理的解釈である。

しかし読者は文献で、他のいくつかの過程が、ときに「デコヒーレンス」と呼ばれたり、形式的、物理的などの意味で、それと「同値」と主張されたりする説明に出会うかもしれない。このような結び付けの大半は、異なる過程がいずれも、系の密度行列の、ある基底における非対角要素の消失に至りうる、という観察に基づく。しかし密度行列による記述は、形式的道具にすぎず、密度行列の形式では効果が似た表現を持つ過程の、決定的な物理的・概念的相違を、やや誤解させやすく、欺くように隠すことを、強調するのが重要である。

% 原文 p.96
\hypertarget{srcpage:96}{}\label{srcpage:96}
たとえば量子工学者や実験家として実際的な立場を採ると、純粋な環境とのもつれとは異なる過程に「デコヒーレンス」という語を結び付けることが、時には実際に正当化される場合もあるが、それらの根本的解釈は、本書の意味での「真の」デコヒーレンスとは異なる。これは単なる哲学的問題ではない。形式的に似た密度行列記述の基礎にある、異なる物理過程を、少なくとも原理上区別できる実験手順が、常に存在する。

Joos\cite{ref017}の作った語を用いた「偽のデコヒーレンス」の、最も一般的で重要な例は、雑音を含む系の異なる実現についての集団平均の結果を、デコヒーレンス過程の記述と解釈することである。これは、密度行列の形式に現れる形式的類似によって、物理的デコヒーレンス過程へ誤って分類する例をなす。この混同の起源は、物理的に非常に異なる二つの設定を記述するために、密度行列を使うことに遡れる。一方で、密度行列の形式は、真の物理的集団、すなわち、各系が純粋状態で記述される$N$個の物理系の集まりを記述するために、しばしば用いられる（2.4.5節の議論も参照）。たとえば各系が、基底状態$\lvert0\rangle$、$\lvert1\rangle$に張られる2次元Hilbert空間で表され、第$i$系の純粋状態が
\begin{equation}\lvert\psi_i\rangle=\frac1{\sqrt2}(\lvert0\rangle+\lvert1\rangle e^{i\phi_i})\tag{2.118}\end{equation}
で与えられるとする。ここで$e^{i\phi_i}$は、各系に固有の相対位相因子である。この位相因子の生じ方は二通りある。第一に、各系が初めから、わずかに異なる初期状態に準備されているかもしれない。あるいは、すべての系が同じ状態から出発しても、わずかに異なるHamiltonianに従うかもしれない。たとえば、各系のHamiltonian$\hat H_i$のポテンシャルが、ランダムな揺らぎ$V_i(t)=V_0(t)+\delta V_i(t)$を含みうる\cite{ref121}。いずれの場合も、異なる系の相対位相因子の違いは、系の準備中、またはその後の発展中に加わる古典的雑音の存在によると考えられる。

$N$個すべての系の純粋状態の集団$\{\lvert\psi_i\rangle\}$について平均すると、集団の密度行列

% 原文 p.97
\hypertarget{srcpage:97}{}\label{srcpage:97}
\begin{equation}\begin{split}\hat\rho&=\frac1N\sum_{i=1}^N\lvert\psi_i\rangle\langle\psi_i\rvert\\&=\frac12\lvert0\rangle\langle0\rvert+\frac12\lvert1\rangle\langle1\rvert\\&\quad+\left(\frac1{2N}\sum_{i=1}^Ne^{i\phi_i}\right)\lvert1\rangle\langle0\rvert+\left(\frac1{2N}\sum_{i=1}^Ne^{-i\phi_i}\right)\lvert0\rangle\langle1\rvert\end{split}\tag{2.119}\end{equation}
を得る。$N$が大きい極限では、ランダムな位相因子$e^{\pm i\phi_i}$の和が、平均して零になるため、
\begin{equation}\hat\rho\longrightarrow\frac12\lvert0\rangle\langle0\rvert+\frac12\lvert1\rangle\langle1\rvert.\tag{2.120}\end{equation}
明らかに、この「集団の位相緩和」は、$\{\lvert0\rangle,\lvert1\rangle\}$基底でデコヒーレンスを起こす個々の系の縮約密度行列と、同じ対角形の密度行列を生む。しかし式(2.119)、(2.120)で干渉項$\lvert0\rangle\langle1\rvert$、$\lvert1\rangle\langle0\rvert$が消えるのは、系の物理的集団の多数の要素について、言い換えれば、特定の雑音過程の異なる実現について、数学的平均を取った結果にすぎない。それは、個々の系の位相コヒーレンスの非局在化を記述する、私たちの意味でのデコヒーレンスには対応しない。今度は密度行列の形式で力学的位相を平均して生じる、「偽のデコヒーレンス」の別の例が、文献\cite{ref122}で批判的に分析されている。

実際、2.4.5節で指摘したように、各々が式(2.118)の純粋状態$\lvert\psi_i\rangle$で記述される、$N$個の物理系の集団の正しい記述は、すべての状態$\lvert\psi_i\rangle$の、純粋状態のテンソル積
\begin{equation}\lvert\Psi\rangle=\prod_{i=1}^N\lvert\psi_i\rangle\tag{2.121}\end{equation}
で与えられる。この状態は$2^N$次元Hilbert空間に「住む」ので、対応する密度行列の大きさは、式(2.119)のような$2\times2$ではなく、$2^N\times2^N$である。環境との相互作用なしでは、各系は完全にユニタリに発展し、したがって位相関係$\phi_i$は、常に明確に保たれる。式(2.119)のように、多くの物理系の集団平均を表すために密度行列の形式を使うと、形式的にだけ位相が「洗い流された」ように見えるが、位相は系から物理的に非局在化してはいないので、少なくとも原理上、各系の観測者には実験的にアクセスできるままである。雑音過程の時間発展を知っていれば、適切なユニタリな逆変換を系に加えて、実験の各実行で、雑音の影響を常に完全に逆転できる。この集団の位相緩和の逆転は、たとえばNMR実験で実際に実装されている。そこでは、雑音により各スピンがわずかに異なる振動数で歳差運動するため、スピンの集まりは、やがて位相が揃わなくなる。しかし、適切に選んだ制御パルスを加えると、スピン集団を「再集束」できる。これはスピンエコー法と呼ばれる。

% 原文 p.98
\hypertarget{srcpage:98}{}\label{srcpage:98}
系のHamiltonianへのランダムな揺らぎの追加である雑音は、系と環境の量子もつれを何も作らず、少なくとも原理上、局所操作によって完全に取り消せる。これに対しデコヒーレンスは、系が環境自由度ともつれ、位相関係が実際上すべての目的に対して不可逆に系から失われ、系への局所測定では、もはやアクセスできなくなることを意味する。それに対応して、系の情報が環境に刻まれる。したがって古典的雑音の場合には、系が環境に乱されるが、式(2.115)は、デコヒーレンスが、系が環境を乱す状況を記述することを示す（文献\cite{ref016}のIV.C節も参照）。すると量子状態の非局所的性格によって、系が環境に与えたこの「歪み」が、今度は系の水準で観測できるものに影響する。これは形式的には縮約密度行列で記述される。しかし2.11節で論じたように、この影響は、環境が与える古典的雑音による系の擾乱よりも、通常ははるかに短い時間尺度で起こる。

確率的揺らぎ、すなわち古典的雑音過程は、環境の系への影響を模擬するためにも、しばしば使われる。その場合、系に加える確率的な「蹴り」によって、デコヒーレンスを模型化する（たとえば文献\cite{ref121,ref123,ref124,ref125}）。もちろん、上で論じたのと同じ問題がここにも当てはまる。古典的雑音を系に加えることは、ランダムに揺らぐHamiltonianで支配される、閉じた系のユニタリ発展に対応する。一方デコヒーレンスは、開いた量子系が、他の量子系、すなわち環境と、決定論的にもつれることで表される。一般に、これら二つの状況から生じる系の量子状態は、異なる観測可能な性質を持つ\cite{ref121}。ある場合にだけ、古典的雑音によるデコヒーレンスの模擬は、密度行列の記述の水準で、形式的に似た結果を与える。しかし、そのような場合でも、雑音とデコヒーレンスの根本的で概念的な相違は残る。7.3節では、量子計算の文脈で、特に、雑音によってデコヒーレンスを実験的・理論的に模擬することに関し、これらの違いをさらに論じる。
